\documentclass[11pt]{article}
\usepackage{amsmath,amssymb,amsthm}

\newtheorem{theorem}{Theorem}
\newtheorem{lemma}[theorem]{Lemma}
\newtheorem{corollary}[theorem]{Corollary}

\title{The $3$-adic collision law for the Collatz terminal decoder}
\author{}
\date{27 August 2026}

\begin{document}
\maketitle

Let
\[
  2^{a_t}x_{t+1}=3x_t+1,
  \qquad
  a_t=v_2(3x_t+1),
  \tag{1}
\]
where every $x_t$ is a positive odd integer.  Put
\[
  Y_t:=x_t+1,
  \qquad
  b_t:=v_3(Y_t),
  \qquad
  Y_t=3^{b_t}u_t,
  \qquad
  3\nmid u_t.
  \tag{2}
\]
For $a\geq1$, define
\[
  r(a):=2^a-2,
  \qquad
  q(a):=v_3(r(a)),
  \tag{3}
\]
where $q(1):=\infty$.

\begin{lemma}[valuation of the decoder forcing term]
For every $a\geq1$,
\[
  q(a)=
  \begin{cases}
    \infty,&a=1,\\
    0,&a\text{ is even},\\
    1+v_3\bigl((a-1)/2\bigr),&a>1\text{ is odd}.
  \end{cases}
  \tag{4}
\]
\end{lemma}

\begin{proof}
The first two cases follow directly from $r(1)=0$ and
$2^a\equiv1\pmod3$ when $a$ is even.  Suppose that $a=2n+1>1$.  Then
\[
  r(a)=2(4^n-1).
  \tag{5}
\]
Write $n=3^s m$, where $3\nmid m$.  The binomial theorem gives
\[
  4^m-1=(1+3)^m-1\equiv3m\pmod9,
\]
so $v_3(4^m-1)=1$.  If $X\equiv1\pmod3$, then
$X^3-1=(X-1)(X^2+X+1)$, while $X^2+X+1\equiv3\pmod9$.  Thus replacing an
exponent by three times that exponent raises the
$3$-adic valuation by exactly one.  Repeating this $s$ times gives
\[
  v_3(4^n-1)=1+s=1+v_3(n),
\]
which proves (4).
\end{proof}

\begin{theorem}[decoder collision law]
Fix one step of (1), and abbreviate
\[
  a:=a_t,
  \qquad
  b:=b_t,
  \qquad
  b':=b_{t+1},
  \qquad
  u:=u_t.
\]
Then
\[
\boxed{
  2^a3^{b'}u_{t+1}=3^{b+1}u+r(a).}
\tag{6}
\]
If $q(a)\neq b+1$, then
\[
\boxed{b'=\min\{q(a),b+1\}.}
\tag{7}
\]
If $q(a)=b+1<\infty$, put
\[
  \rho_a:=\frac{r(a)}{3^{q(a)}}.
  \tag{8}
\]
Then $3\nmid\rho_a$ and
\[
\boxed{
  b'=b+1+v_3(u+\rho_a).}
\tag{9}
\]
In particular,
\[
  a=1\Longrightarrow b'=b+1,
  \qquad
  a\equiv0\pmod2\Longrightarrow b'=0.
  \tag{10}
\]
\end{theorem}

\begin{proof}
Adding $2^a-1$ to both sides of (1) gives
\[
  2^a(x_{t+1}+1)=3(x_t+1)+2^a-2.
\]
Substitution from (2) proves (6).  Since $2^a$ is a unit modulo $3$, the
number $b'$ is the valuation of the right side of (6).

The two summands on that side have valuations $b+1$ and $q(a)$.  If these
valuations are distinct, the valuation of their sum is their minimum.  This
proves (7).  If they are equal and finite, factoring out $3^{b+1}$ gives
\[
  3^{b+1}(u+\rho_a),
\]
which proves (9).  The two statements in (10) follow from (4), (7), and
$q(1)=\infty$.
\end{proof}

\begin{corollary}[bounded-ceiling exponent sieve]
Fix $D\geq1$, and suppose that $0\leq b\leq D$.  The condition $b'\leq D$
is characterized as follows.

\begin{enumerate}
  \item If $b=D$, then
  \[
  \boxed{
    b'\leq D
    \quad\Longleftrightarrow\quad
    a\not\equiv1\pmod{2\cdot3^D}.}
  \tag{11}
  \]

  \item If $b<D$ and $q(a)=b+1$, then
  \[
  \boxed{
    b'\leq D
    \quad\Longleftrightarrow\quad
    u\not\equiv-\rho_a\pmod{3^{D-b}}.}
  \tag{12}
  \]

  \item If $b<D$ and $q(a)\neq b+1$, then $b'\leq D$ automatically.
\end{enumerate}

Consequently, if $b_t\leq D$ on a tail, then every time with $b_t=D$
satisfies
\[
  a_t\not\equiv1\pmod{2\cdot3^D},
  \tag{13}
\]
and that tail contains no factor $1^{D+1}$.
\end{corollary}

\begin{proof}
Suppose first that $b=D$.  If $q(a)\geq D+1$, equations (7) and (9) give
$b'\geq D+1$.  If $q(a)\leq D$, equation (7) gives $b'=q(a)\leq D$.
Therefore,
\[
  b'\leq D\quad\Longleftrightarrow\quad q(a)\leq D.
  \tag{14}
\]
By (4), the complementary condition $q(a)\geq D+1$ is equivalent to
\[
  a\equiv1\pmod{2\cdot3^D},
\]
including the case $a=1$.  This proves (11).

Now suppose that $b<D$.  If $q(a)=b+1$, equation (9) gives
\[
  b'\leq D
  \quad\Longleftrightarrow\quad
  v_3(u+\rho_a)\leq D-b-1.
\]
This is exactly (12).  If $q(a)<b+1$, then (7) gives $b'=q(a)\leq b<D$.
If $q(a)>b+1$, the same equation gives $b'=b+1\leq D$.  This proves the
third assertion.

Finally, equation (13) is (11) applied along the tail.  Each exponent-one
step raises the decoder depth by one according to (10).  A factor
$1^{D+1}$ would therefore raise a nonnegative decoder depth above $D$,
which is impossible.
\end{proof}

\begin{theorem}[one-forbidden-class decoder automaton]
Fix $D\geq0$, and put
\[
  M_D:=3^{D+1},
  \qquad
  H_D:=2\cdot3^D.
  \tag{15}
\]
For $s\in\mathbb Z/M_D\mathbb Z$ and $a\geq1$, define
\[
  F_a(s):=2^{-a}\bigl(3s+2^a-2\bigr)
  \pmod {M_D}.
  \tag{16}
\]
If $s_t\equiv x_t+1\pmod {M_D}$, then
\[
  s_{t+1}=F_{a_t}(s_t).
  \tag{17}
\]
Moreover, for every nonzero state $s$, there is exactly one residue class
$\chi_D(s)\pmod {H_D}$ such that
\[
\boxed{
  F_a(s)=0
  \quad\Longleftrightarrow\quad
  a\equiv\chi_D(s)\pmod {H_D}.}
\tag{18}
\]
Equivalently, at every state with decoder depth at most $D$, exactly one
exponent class modulo $2\cdot3^D$ raises the next decoder depth above $D$.
\end{theorem}

\begin{proof}
Equation (17) is (6) reduced modulo $M_D$.  We next verify that $2$ generates
the unit group modulo $3^k$ for every $k\geq1$.  The calculation in the proof
of Lemma~1 gives
\[
  v_3(2^{2m}-1)=v_3(4^m-1)=1+v_3(m).
  \tag{19}
\]
An exponent $h$ satisfying $2^h\equiv1\pmod {3^k}$ must be even, say
$h=2m$, and (19) shows that its least positive value is
$2\cdot3^{k-1}$.  This equals $\varphi(3^k)$, so $2$ is a generator.

Now $F_a(s)=0$ is equivalent to
\[
  2^a\equiv2-3s\pmod {M_D}.
  \tag{20}
\]
The right side is a unit, and the powers of $2$ run once through all units as
$a$ runs through a complete residue system modulo $H_D$.  Hence (20) selects
exactly one exponent class, proving (18).
\end{proof}

\begin{theorem}[exact finite realization of decoder paths]
Fix $D\geq0$ and a nonzero initial state
$s_0\in\mathbb Z/M_D\mathbb Z$.  Let
$(a_0,\ldots,a_{L-1})$ be any positive exponent word, and define
recursively
\[
  s_{t+1}:=F_{a_t}(s_t).
  \tag{21}
\]
If $s_t\neq0$ for $0\leq t\leq L$, then infinitely many positive odd
integers $n$ have this exact accelerated exponent prefix and satisfy
\[
  x_t+1\equiv s_t\pmod {M_D}
  \qquad(0\leq t\leq L).
  \tag{22}
\]
In particular, their decoder depths obey $v_3(x_t+1)\leq D$ throughout the
displayed finite orbit segment.

If $D\geq1$ and the exponent alphabet is
\[
  \mathcal A_D:=\{1,2,\ldots,H_D\},
  \tag{23}
\]
then exactly $(H_D-1)^L$ words in $\mathcal A_D^L$ generate a nonzero path
from any prescribed nonzero state $s_0$.  Every one of those words has the
positive realizations asserted above.
\end{theorem}

\begin{proof}
We first recall the finite cylinder fact and include its induction.  A
positive exponent word of total exponent $A$ is realized exactly by one odd
residue class modulo $2^{A+1}$.  The empty word gives the unique odd class
modulo $2$.  Suppose that a prefix of length $t$ and total exponent $A$ has
representative $\rho$.  Replacing its initial value by
$\rho+\ell2^{A+1}$ replaces its endpoint by
\[
  x_t(\rho)+2\ell3^t.
  \tag{24}
\]
To append the exponent $a$, exactness is equivalent to
\[
  3x_t(\rho)+1+6\ell3^t\equiv2^a\pmod {2^{a+1}}.
  \tag{25}
\]
After division by $2$, the coefficient of $\ell$ is the odd number
$3^{t+1}$.  Thus (25) selects one class $\ell\pmod {2^a}$, completing the
induction.

Let $\rho$ be the resulting cylinder class for the displayed word and put
$A:=\sum_{t=0}^{L-1}a_t$.  The simultaneous congruences
\[
  n\equiv\rho\pmod {2^{A+1}},
  \qquad
  n\equiv s_0-1\pmod {M_D}
  \tag{26}
\]
have one class modulo $2^{A+1}M_D$ by the Chinese remainder theorem.  This
class contains infinitely many positive odd integers.  Its first congruence
gives the exact exponent word.  Its second congruence, followed by (17),
proves (22) by induction.

Finally, the alphabet (23) contains one representative of every exponent
class modulo $H_D$.  Theorem~4 says that exactly one of these $H_D$ letters
leaves the nonzero state space at each step.  Hence the rooted path tree has
exactly $H_D-1$ children at every vertex and exactly $(H_D-1)^L$ paths of
length $L$.  The first part of the theorem realizes each path.
\end{proof}

\end{document}
